\chapter{Stabile}

\section{Stabilized parameterizations of $K$, $H$, and $R$}

In the baseline specification, the latent-factor dynamics are parameterized by three neural
networks: a drift network $K$ of the form $\mu(z)=Mz+N$, a volatility network $H$
mapping the latent state to $(\log \sigma_1,\log \sigma_2,\operatorname{atanh}(\rho))$,
and a short-rate network $R$ mapping the latent state to $\tilde r(z)$. In the present
two-factor setting, these networks determine the coefficients of the risk-neutral dynamics
in \eqref{eq:30}--\eqref{eq:33}. While this baseline parameterization is flexible, we found
that it may lead to numerical instability when the model is simulated forward or when the
latent state moves outside regions well covered by the training data. We therefore replace
the original forms of $K$, $H$, and $R$ by stabilized parameterizations that preserve the
financial interpretation of the model while enforcing basic structural regularity.

\paragraph{Stabilization of the drift network $K$.}
In the baseline model, the drift is specified as
\begin{equation}
\mu(z)=Mz+N,
\label{eq:k_baseline_combined}
\end{equation}
where $M \in \mathbb{R}^{d\times d}$ and $N \in \mathbb{R}^d$ are trainable parameters.
Although this form is intended to model mean reversion, it does not guarantee that the
learned matrix $M$ is stable. In particular, if $M$ has an eigenvalue with non-negative
real part, the deterministic part of the latent dynamics may fail to be mean-reverting.

To enforce stability by construction, we replace $M$ by the parameterization
\begin{equation}
M = -\left(V^\top V + \varepsilon I_d\right),
\label{eq:k_stable_combined}
\end{equation}
where $V \in \mathbb{R}^{d\times d}$ is trainable, $I_d$ is the identity matrix, and
$\varepsilon>0$ is a small constant. The stabilized drift is therefore
\begin{equation}
\mu(z)= -\left(V^\top V + \varepsilon I_d\right)z + N.
\label{eq:k_stable_drift_combined}
\end{equation}
Since $V^\top V$ is positive semidefinite, the matrix $V^\top V+\varepsilon I_d$ is
strictly positive definite, and hence $M$ is negative definite. Indeed, for every nonzero
$x \in \mathbb{R}^d$,
\begin{equation}
x^\top Mx
=
-x^\top\left(V^\top V+\varepsilon I_d\right)x
=
-\|Vx\|^2-\varepsilon\|x\|^2
<0.
\end{equation}
It follows that all eigenvalues of $M$ are real and strictly negative. The linear drift is
therefore mean-reverting by construction. The corresponding equilibrium point is
\begin{equation}
z^\ast = \left(V^\top V+\varepsilon I_d\right)^{-1}N,
\end{equation}
so that the drift may equivalently be written as
\begin{equation}
\mu(z)=-Q(z-z^\ast),
\qquad
Q:=V^\top V+\varepsilon I_d.
\end{equation}
This preserves the original linear structure of $K$, while ensuring a stable latent drift.

\paragraph{Stabilization of the volatility network $H$.}
In the baseline model, the network $H$ outputs transformed volatility parameters,
\begin{equation}
H(z)=
\begin{pmatrix}
\log \sigma_1(z) \\
\log \sigma_2(z) \\
\operatorname{atanh}(\rho(z))
\end{pmatrix},
\label{eq:h_baseline_combined}
\end{equation}
which are then mapped back through exponentiation and the hyperbolic tangent. This
ensures $\sigma_i(z)>0$ and $\rho(z)\in(-1,1)$, but it does not prevent the volatilities
from becoming arbitrarily small or arbitrarily large, nor does it prevent the correlation
from approaching $\pm 1$ too closely. In the two-factor diffusion matrix
\begin{equation}
\sigma(z)=
\begin{pmatrix}
\sigma_1(z) & 0 \\
\rho(z)\sigma_2(z) & \sqrt{1-\rho(z)^2}\,\sigma_2(z)
\end{pmatrix},
\label{eq:h_diffusion_combined}
\end{equation}
such behavior may lead to poor numerical conditioning.

To regularize these outputs, we replace the original transformation by the smooth bounded
parameterization
\begin{equation}
\log \sigma_i(z)
=
\log \sigma_{\min}
+
\bigl(\log \sigma_{\max}-\log \sigma_{\min}\bigr)
\,\operatorname{sigmoid}\!\bigl(h_i(z)+c_i\bigr),
\qquad i=1,2,
\label{eq:h_stable_sigma_combined}
\end{equation}
and
\begin{equation}
\rho(z)=\rho_{\max}\tanh\!\bigl(h_\rho(z)\bigr),
\label{eq:h_stable_rho_combined}
\end{equation}
where $h_1(z)$, $h_2(z)$ and $h_\rho(z)$ denote the raw outputs of the final layer,
$0<\sigma_{\min}<\sigma_{\max}$ are fixed bounds, $\rho_{\max}\in(0,1)$ is fixed, and
$c_1,c_2$ are learnable offsets. Consequently,
\begin{equation}
\sigma_i(z)\in(\sigma_{\min},\sigma_{\max}),
\qquad i=1,2,
\end{equation}
and
\begin{equation}
\rho(z)\in(-\rho_{\max},\rho_{\max}),
\qquad \rho_{\max}<1.
\end{equation}
Hence, the volatilities remain strictly positive and uniformly bounded, while the correlation is kept strictly away from the singular values $\pm 1$. In particular, the term $\sqrt{1-\rho(z)^2}$ in \eqref{eq:h_diffusion_combined} remains real and bounded away from zero, which improves the conditioning of the diffusion matrix during both training and simulation.

\paragraph{Stabilization of the short-rate network $R$.}
In the baseline model, the short rate is represented by a feedforward network
\begin{equation}
\tilde r(z)=R_\theta(z),
\label{eq:r_baseline_combined}
\end{equation}
with unconstrained scalar output. While flexible, this form allows the network to produce implausibly large positive or negative short rates when evaluated outside the region well covered by the training data. Since $\tilde r(z)$ enters directly into the no-arbitrage ODE system and the bank-account process, such values may create numerical instability.

We therefore replace the unconstrained output layer by the bounded parameterization
\begin{equation}
\tilde r(z)=r_{\mathrm{center}} + r_{\mathrm{scale}}\tanh\!\bigl(f_\theta(z)\bigr),
\label{eq:r_stable_combined}
\end{equation}
where $f_\theta:\mathbb{R}^d\to\mathbb{R}$ has the same hidden architecture as the original network, $r_{\mathrm{center}}\in\mathbb{R}$ is a learnable center, and $r_{\mathrm{scale}}>0$ is a learnable scale. To enforce positivity of the scale, we write
\begin{equation}
r_{\mathrm{scale}}=\operatorname{softplus}(a)+\varepsilon_r,
\label{eq:r_scale_combined}
\end{equation}
where $a\in\mathbb{R}$ is unconstrained and $\varepsilon_r>0$ is a small constant.
Since $\tanh(x)\in(-1,1)$ for all $x$, the stabilized short rate satisfies
\begin{equation}
\tilde r(z)\in
\left(
r_{\mathrm{center}}-r_{\mathrm{scale}},
\;
r_{\mathrm{center}}+r_{\mathrm{scale}}
\right).
\end{equation}
Thus, the short rate remains in a learnable but controlled interval. This is a smooth reparameterization rather than a hard clipping rule, and it therefore preserves differentiability while preventing extreme outputs.

\paragraph{Summary.}
The three modifications above target different sources of instability. The stabilized
drift network $K$ guarantees a stable mean-reverting linear drift; the stabilized
volatility network $H$ keeps volatilities bounded and correlations away from $\pm 1$;
and the stabilized short-rate network $R$ prevents extreme short-rate realizations.
Together, these changes preserve the economic interpretation of the original
parameterization while improving the numerical robustness of the latent-factor dynamics
used in the arbitrage-free decoder.


\chapter{Numerical Implementation and Estimation Strategy}
\label{chap:numerical_implementation}

This section describes how the affine-inspired encoder--decoder framework developed in Section~4 is implemented numerically and trained on the data described in Section~5. The implementation follows the theoretical construction closely: the encoder maps observed swap rates to a latent state, the decoder produces the quantities entering the no-arbitrage ODE system, and the resulting bond prices are transformed into model-implied swap rates. Since the arbitrage restrictions are embedded directly in the bond-pricing structure, they remain satisfied throughout.

The model is implemented in Python using the PyTorch library, which provides efficient tensor-based computation and automatic differentiation, allowing the entire pricing pipeline, including the numerical ODE solver, to be treated as a differentiable computational graph. In the following, we set $M=8$ corresponding to the maturity grid
\[
\{1\text{Y},2\text{Y},3\text{Y},5\text{Y},10\text{Y},15\text{Y},20\text{Y},30\text{Y}\},
\]
and consider latent dimensions $l \in \{1,2,3,4\}$.

\section{Forward Pass and Numerical Solution}
\label{sec:forward_pass_numerical_solution}

For each observation date $t$, the input consists of the observed par swap rate vector $\mathbf{S}_t \in \mathbb{R}^8$. The forward pass mirrors the theoretical architecture introduced in Section~\ref{sec:theory_placeholder}. % TODO: replace with correct reference
First, the observed swap curve is encoded into a latent state
\[
z_t = A_1 \mathbf{S}_t,
\qquad
z_t \in \mathbb{R}^l,
\]
where $A_1 \in \mathbb{R}^{l \times 8}$ is the learned encoder matrix.

Conditional on $z_t$, the decoder evaluates $\mathcal{G}(z_t,\tau)$ together with the parameterizations of the drift $\mu(z_t)$, the volatility matrix $\sigma(z_t)$, and the short rate $\tilde r(z_t)$. The drift is modeled as an affine transformation of the latent state, while the volatility and short-rate functions are obtained from the neural networks $\mathcal{H}$ and $\mathcal{R}$, respectively. The volatility network output is constructed so that the instantaneous covariance matrix
\[
\Sigma(z_t) = \sigma(z_t)\sigma(z_t)^\top
\]
is positive semidefinite by construction.

Using automatic differentiation, the derivatives of $\mathcal{G}(z_t,\tau)$ required for the coefficient functions $\alpha_{z_t}(\tau)$, $\beta_{z_t}(\tau)$, and $\gamma_{z_t}(\tau)$ are computed on
\[
\tau \in \{0,1,\dots,30\}.
\]
The inclusion of $\tau=0$ allows the boundary conditions $A_{z_t}(0)=B_{z_t}(0)=0$ to be imposed directly.

For each realized latent state $z_t$, the coupled ODE system
\begin{align}
\partial_\tau B_{z_t}(\tau) &= \alpha_{z_t}(\tau) B_{z_t}(\tau) + \beta_{z_t}(\tau), \label{eq:B_ode_num} \\
\partial_\tau A_{z_t}(\tau) &= B_{z_t}(\tau)^2 \gamma_{z_t}(\tau), \label{eq:A_ode_num}
\end{align}
is solved numerically on this grid using a fourth-order Runge--Kutta scheme with the $3/8$-rule. % TODO: add appendix reference if relevant

The resulting functions $A_{z_t}(\tau)$ and $B_{z_t}(\tau)$ are substituted into (4.3), yielding zero-coupon bond prices on maturities $\tau \in \{1,\dots,30\}$. Model-implied swap rates are then obtained from these bond prices using (4.1), evaluated at the $M=8$ observed maturities.

Each forward pass therefore maps an observed swap curve $\mathbf{S}_t$ to a reconstructed swap curve $\widetilde{\mathbf{S}}_t$ through the sequence
\[
\mathbf{S}_t
\longrightarrow
z_t
\longrightarrow
(\mathcal{G},\mu,\sigma,\tilde r)
\longrightarrow
(\alpha,\beta,\gamma)
\longrightarrow
(A,B)
\longrightarrow
P(z_t,\tau)
\longrightarrow
\widetilde{\mathbf{S}}_t.
\]

Because all steps are implemented using differentiable tensor operations, the entire pricing map remains differentiable with respect to the model parameters.

\section{Numerical Arbitrage Diagnostics}
\label{sec:numerical_arbitrage_diagnostics}

Although the decoder is constructed so that the no-arbitrage restriction is satisfied by design, it is still useful to verify this property numerically after solving the ODE system. To this end, we compute a maturity-by-maturity residual of the bond-pricing equation together with an associated approximate Sharpe-ratio diagnostic.

For each observation date $t$ and each maturity $\tau$ on the grid, the model produces the quantities
\[
\mathcal{G}(z_t,\tau),
\qquad
A_{z_t}(\tau),
\qquad
B_{z_t}(\tau),
\]
as well as the risk-neutral objects
\[
\mu(z_t),
\qquad
\sigma(z_t),
\qquad
\tilde r(z_t).
\]

Using automatic differentiation, we also compute the maturity derivative $\partial_\tau \mathcal{G}(z_t,\tau)$, the gradient $\nabla_z \mathcal{G}(z_t,\tau)$, and the trace term
\[
\mathrm{Tr}\!\left(\sigma(z_t)^\top H_z \mathcal{G}(z_t,\tau)\sigma(z_t)\right),
\]
where $H_z \mathcal{G}$ denotes the Hessian of $\mathcal{G}$ with respect to the latent state.

The numerical ODE solver delivers the functions $A_{z_t}(\tau)$ and $B_{z_t}(\tau)$ on the maturity grid. Their derivatives with respect to maturity are obtained directly from the ODE system,
\begin{align}
\partial_\tau B_{z_t}(\tau) &= \alpha_{z_t}(\tau) B_{z_t}(\tau) + \beta_{z_t}(\tau), \label{eq:B_ode_diag} \\
\partial_\tau A_{z_t}(\tau) &= \gamma_{z_t}(\tau) B_{z_t}(\tau)^2. \label{eq:A_ode_diag}
\end{align}

Combining these objects, we define the bond-pricing residual
\begin{equation}
R_t(\tau)
=
- \tilde r(z_t)
- \partial_\tau A_{z_t}(\tau)
+ \mathcal{G}(z_t,\tau)\partial_\tau B_{z_t}(\tau)
+ B_{z_t}(\tau)\Xi_t(\tau)
+ B_{z_t}(\tau)^2 \gamma_{z_t}(\tau),
\label{eq:bond_pricing_residual}
\end{equation}
where
\begin{equation}
\Xi_t(\tau)
=
\partial_\tau \mathcal{G}(z_t,\tau)
-
\nabla_z \mathcal{G}(z_t,\tau)^\top \mu(z_t)
-
\frac{1}{2}
\mathrm{Tr}\!\left(\sigma(z_t)^\top H_z \mathcal{G}(z_t,\tau)\sigma(z_t)\right).
\label{eq:Xi_term}
\end{equation}

If the decoded bond prices satisfy the no-arbitrage pricing relation exactly, then $R_t(\tau)=0$ for all maturities. In practice, small nonzero values may arise because of numerical approximation, finite solver tolerance, and automatic-differentiation error.

To obtain a scale-normalized diagnostic comparable across maturities, we follow the Sharpe-ratio logic discussed in the theory chapter and define the approximate Sharpe ratio
\begin{equation}
\mathrm{SR}_t(\tau) = \frac{R_t(\tau)}{\tau \bar{\sigma}},
\label{eq:approx_sharpe_ratio}
\end{equation}
where $\bar{\sigma} > 0$ is a fixed volatility scaling constant. In the implementation, we set
\[
\bar{\sigma} = 0.006.
\]

The quantities $R_t(\tau)$ and $\mathrm{SR}_t(\tau)$ are computed on the full maturity grid $\tau = 1,\dots,30$ for each observation. They are not included in the training loss, but are retained as post-estimation diagnostics of the numerical consistency of the decoded bond curve. In the empirical analysis, we summarize these diagnostics across dates, maturities, and currencies, and also report representative maturity profiles.

\section{Estimation Objective}
\label{sec:estimation_objective}

Let $\Omega$ denote the collection of all trainable parameters in the model, including the encoder matrix $A_1$ and the parameters governing the decoder components $\mathcal{G}$, $\mathcal{K}$, $\mathcal{H}$, and $\mathcal{R}$. The model is trained by minimizing the mean squared error between observed and reconstructed swap rates across the sample,
\begin{equation}
\mathcal{L}(\Omega)
=
\frac{1}{N}
\sum_{t=1}^{N}
\sum_{\tau \in M}
\left(
S_t(\tau) - \widetilde{S}_t(\tau;\Omega)
\right)^2,
\label{eq:training_loss}
\end{equation}
where $N$ denotes the number of observations. The estimator is defined as
\begin{equation}
\hat{\Omega}
=
\arg\min_{\Omega}\mathcal{L}(\Omega).
\label{eq:estimator}
\end{equation}

\section{Gradient-Based Optimization}
\label{sec:gradient_based_optimization}

The optimization problem in \eqref{eq:training_loss} is solved using stochastic gradient-based methods, as described in Section~\ref{sec:ml_placeholder}. % TODO: replace with correct machine-learning reference
Parameter updates are obtained by backpropagating the loss through the full computational graph.

A key feature of the implementation is that differentiation extends through the numerical ODE solver.

Since the Runge--Kutta recursion is implemented using elementary tensor operations, automatic differentiation propagates gradients through
\begin{align*}
\text{Step 1} \quad & \text{the encoder mapping } A_1\mathbf{S}_t, \\
\text{Step 2} \quad & \text{the parameterizations of } \mathcal{G}(z_t,\tau),\ \mu(z_t),\ \sigma(z_t),\ \text{and } \tilde r(z_t), \\
\text{Step 3} \quad & \text{the construction of the coefficient functions } \alpha_{z_t}(\tau),\ \beta_{z_t}(\tau),\ \text{and } \gamma_{z_t}(\tau), \\
\text{Step 4} \quad & \text{the numerical solution of the ODE system for } A_{z_t}(\tau)\ \text{and } B_{z_t}(\tau), \\
\text{Step 5} \quad & \text{and the final transformation from bond prices to swap rates.}
\end{align*}

Each parameter update therefore reflects the effect of the corresponding model component on the entire reconstructed term structure.

Training is performed using mini-batch stochastic gradient descent together with the Adam optimizer and learning rate scheduling, as discussed in the machine-learning section. % TODO: if already stated elsewhere, shorten this sentence there or here

\paragraph{Conclusion.}
The fitted parameters $\hat{\Omega}$ define a term structure model that simultaneously extracts a low-dimensional representation of the observed swap curve and prices zero-coupon bonds consistently with the risk-neutral no-arbitrage condition.